%--------------------------
\begin{frame}{Question 5}

\begin{block}{}
La date de dépôt d’une demande nationale antérieure sera reconnue aux fins de la 
phase internationale comme date de priorité dans une demande PCT :

\begin{itemize}
  \item[$\bullet$] a. si la priorité de ladite demande antérieure est valablement revendiquée dans la 
  demande PCT et qu’elle respecte les exigences de la Convention de Paris. 
  \item[$\circ$] b. même si la date de dépôt de cette demande nationale est antérieure de plus de 12 
  mois à la date de dépôt PCT, et qu’elle respecte les autres exigences de la 
  Convention de Paris. 
  \item[$\bullet$] c. même si la date de dépôt de cette demande nationale est antérieure de plus de 12 
  mois à la date du dépôt PCT, mais à condition qu’elle tombe dans le délai de 12 + 2 
  mois avant ladite date de dépôt PCT, et qu’elle respecte les autres exigences de la 
  Convention de Paris. 
  \item[$\circ$] d. Aucune des réponses a, b ou c n’est correcte. 
\end{itemize}

\end{block}

\end{frame}

% %--------------------------
% \begin{frame}{Question 5 - à revoir}

% \begin{itemize}
%   \item[$\bullet$] \textbf{a.} \; Correct : une priorité valablement revendiquée dans les 12 mois
%   produit effet conformément à la Convention de Paris
%   (\pctart{8}(1) + \pctart{8}(2)(a)).

%   \item[$\circ$] \textbf{b.} \; Faux : une priorité revendiquée après l’expiration des 12 mois
%   n’est pas reconnue automatiquement.

%   \item[$\bullet$] \textbf{c.} \; Correct : une restauration du droit de priorité est possible
%   si la demande PCT est déposée dans un délai de 2 mois après l’expiration des 12 mois
%   (\pctrule{26bis}.3).

%   \item[$\circ$] \textbf{d.} \; Faux puisque \textbf{a} et \textbf{c} sont correctes.
% \end{itemize}

% \end{frame}
